\begin{proof}[Proof of \cref{obj:thm:sharp-local-linear-constant-and-representers}]
Throughout, when \(t\) is fixed we write \(n=n_t\), \(d=d_t\), \(m=m_t\), \(A=A_{d_t}\) and \(\bar\beta=\bar\beta_{d_t}=\min\{\beta,d\}\).  The population \(V_{n}\) is partitioned into the blocks \(V_1,\ldots,V_m\), each of size \(d\); \(b(i)\) is the block containing \(i\), \(Z_b=(Z_j)_{j\in V_b}\), and \(Z_S=\prod_{j\in S}Z_j\).  Since \(G_t\) is the directed complete-block graph,
\[
N_i^{G_t}=V_{b(i)}\quad\text{for every }i\in V_{n},
\qquad
n=md .
\]
The clauses are proved in the order in which their ingredients are established: the exact worst-case formula, then the per-block lower bound, then the minimax constant, then the two asymptotic criteria, then the witness sequence.

We first record the three block-score identities that every step uses.  The block score \(g_{d}\) of \cref{obj:def:block-score-energy} lies in \(\mathcal P_{d}\), and \cref{obj:lem:block-energy-representer} gives, under \cref{obj:ass:bernoulli-design},
\[
\mathbb E_Z[g_{d}(Z_b)]=0,
\qquad
\mathbb E_Z[g_{d}(Z_b)Z_S]=1
\quad\text{for every }\varnothing\ne S\subseteq V_b\text{ with }|S|\le\bar\beta ,
\]
\[
\mathbb E_Z[g_{d}(Z_b)^2]=L_{d}(g_{d})=A>0 .
\]
The first two displays are the representer identity \(\mathbb E_Z[g_{d}(Z)f(Z)]=L_{d}(f)\) at \(f\equiv1\) and at \(f=Z_S\); the third combines that identity at \(f=g_{d}\) with \(\mathbb E_Z[h_{d}(Z)^2]=A^{-1}\) and \(g_{d}=Ah_{d}\), and the strict positivity is the energy-scale bound \(A\ge c_1\binom{d}{k_\star(d,\beta,p)}>0\). % lean: blockScore_mean_zero, blockScore_raw_moment, snipeScore_sq_expectation_global, blockEnergy_pos

\begin{enumerate}
\item \textbf{The exact worst-case risk formula.}  Fix \(t\) and \(w\in\mathcal E_t^{\mathrm{loc,lin}}\).  For \(c\in\mathcal M_t(G_t)\) the potential outcomes are \(Y_i(Z)=\sum_{S\subseteq V_{b(i)},\,|S|\le\bar\beta}c_{i,S}Z_S\), and the target is the all-treated versus all-control contrast
\[
\tau_{n}
=
\frac1n\sum_{i\in V_{n}}\{Y_i(\mathbf 1_{n})-Y_i(\mathbf 0_{n})\}
=
\frac1n\sum_{i\in V_{n}}
\sum_{\substack{S\subseteq V_{b(i)}\\ |S|\le\bar\beta}}
c_{i,S}\mathbf 1(S\ne\varnothing).
\]
Subtracting this from \(\widehat\tau_{w,t}=n^{-1}\sum_i w_{i,t}(Z_{b(i)})Y_i(Z)\) and grouping units by block,
\[
\widehat\tau_{w,t}-\tau_{n}
=
\frac1n\sum_{b=1}^{m}e_{b}(c;Z),
\qquad
e_{b}(c;Z)
:=
\sum_{i\in V_b}
\sum_{\substack{S\subseteq V_b\\ |S|\le\bar\beta}}
c_{i,S}\bigl\{w_{i,t}(Z_b)Z_S-\mathbf 1(S\ne\varnothing)\bigr\}.
\]
% lean: locLinEstimator_error_expansion, sum_blockScheduleError_eq_total
Each \(e_{b}(c;Z)\) is a function of \(Z_b\) alone, and the moment restrictions defining \(\mathcal E_t^{\mathrm{loc,lin}}\) give \(\mathbb E_Z[e_{b}(c;Z)]=0\) termwise.  Distinct blocks use disjoint assignment coordinates, so under \cref{obj:ass:bernoulli-design} the variables \(e_1(c;Z),\ldots,e_{m}(c;Z)\) are independent and the cross terms vanish:
\[
\mathbb E_Z\!\left[(\widehat\tau_{w,t}-\tau_{n})^2\right]
=
\frac{1}{n^2}\sum_{b=1}^{m}\mathbb E_Z\!\left[e_{b}(c;Z)^2\right].
\]
% lean: E_completeBlock_scheduleError_sum_sq, locLinRiskAt_block_decomposition
For a fixed block, \(c\mapsto e_{b}(c;Z)\) is linear in the rows \((c_{i,S})_{S}\) and the only constraint \(\mathcal M_t(G_t)\) places on a row is \(\sum_{S}|c_{i,S}|\le B\) with \(S\subseteq V_b\) and \(|S|\le\bar\beta\).  The criterion \(c\mapsto\mathbb E_Z[e_{b}(c;Z)^2]\) is convex in each row, so its maximum over the row simplex of radius \(B\) is attained at a vertex \(c_{i,S}=\varepsilon_iB\mathbf 1(S=S_i)\) with \(\varepsilon_i\in\{-1,1\}\), for which \(e_{b}(c;Z)=B\sum_{i\in V_b}\varepsilon_i\{w_{i,t}(Z_b)Z_{S_i}-\mathbf 1(S_i\ne\varnothing)\}\).  Hence
\[
\sup_{c\in\mathcal M_t(G_t)}\mathbb E_Z\!\left[e_{b}(c;Z)^2\right]
=
B^2\Psi_{b,t}(w).
\]
% lean: blockScheduleError_sq_le_extremal, exists_blockExtremal_choice, extremeBlockSchedule
The row constraint is unitwise, so the maximizing vertices may be chosen block by block and assembled into one schedule in \(\mathcal M_t(G_t)\) that attains every block maximum simultaneously.  Therefore
\[
\sup_{c\in\mathcal M_t(G_t)}
\mathbb E_Z\!\left[(\widehat\tau_{w,t}-\tau_{n})^2\right]
=
\frac{B^2}{n^2}\sum_{b=1}^{m}\Psi_{b,t}(w),
\]
which is the asserted worst-case identity. % lean: locLinWorstRisk_exact_blockExtremal

\item \textbf{The per-block lower bound.}  Fix \(t\), \(w\in\mathcal E_t^{\mathrm{loc,lin}}\) and a block \(V_b\).  Every \(i\in V_b\) has \(N_i^{G_t}=V_b\), so the single function \(g_{d}(Z_b)\) is the common block score for all of them.  Expand it in the raw monomials it is spanned by,
\[
g_{d}(Z_b)=\sum_{\substack{T\subseteq V_b\\ |T|\le\bar\beta}}a_TZ_T,
\qquad
\sum_{\substack{T\subseteq V_b\\ |T|\le\bar\beta}}a_T=g_{d}(\mathbf 1_{d}),
\qquad
a_\varnothing=g_{d}(\mathbf 0_{d}).
\]
Applying \(\mathbb E_Z[w_{i,t}]=0\) to the term \(T=\varnothing\) and \(\mathbb E_Z[w_{i,t}Z_T]=1\) to every other term gives, for every \(i\in V_b\),
\[
\mathbb E_Z[w_{i,t}(Z_b)g_{d}(Z_b)]
=
\sum_{\substack{\varnothing\ne T\subseteq V_b\\ |T|\le\bar\beta}}a_T
=
g_{d}(\mathbf 1_{d})-g_{d}(\mathbf 0_{d})
=
L_{d}(g_{d})
=
A .
\]
% lean: locLinWeight_snipeScore_pair
Summing over the \(d\) units of \(V_b\) and using \(\mathbb E_Z[g_{d}(Z_b)^2]=A\), the expansion of a nonnegative square gives
\[
0\le
\mathbb E_Z\!\left[
\left\{\sum_{i\in V_b}w_{i,t}(Z_b)-dg_{d}(Z_b)\right\}^2
\right]
=
\mathbb E_Z\!\left[
\left\{\sum_{i\in V_b}w_{i,t}(Z_b)\right\}^2
\right]
-2d(dA)+d^2A ,
\]
that is,
\[
d^2A
\le
\mathbb E_Z\!\left[
\left\{\sum_{i\in V_b}w_{i,t}(Z_b)\right\}^2
\right].
\]
The choice \(\varepsilon_i=1\) and \(S_i=\varnothing\) for every \(i\in V_b\) is admissible in \(\Psi_{b,t}\), since \(\varnothing\subseteq V_b\) and \(|\varnothing|=0\le\bar\beta\), and it produces exactly that expectation; as \(\Psi_{b,t}\) is a maximum over admissible choices,
\[
\mathbb E_Z\!\left[
\left\{\sum_{i\in V_b}w_{i,t}(Z_b)\right\}^2
\right]
\le\Psi_{b,t}(w),
\qquad\text{hence}\qquad
d_t^2A_{d_t}\le \Psi_{b,t}(w).
\]
% lean: blockExtremal_all_empty_le, completeBlock_blockExtremal_lower

\item \textbf{The exact minimax constant.}  Define the canonical block-local weights
\[
w^0_{i,t}(Z):=g_{d}(Z_{b(i)}),
\]
which are the SNIPE scores of \cref{obj:def:snipe-score} on \(G_t\) because \(N_i^{G_t}=V_{b(i)}\) has exactly \(d\) elements.  They depend on \(Z\) only through \(Z_{b(i)}\), and the three block-score identities recorded above are precisely the centering and unit raw-moment restrictions, so \(w^0\in\mathcal E_t^{\mathrm{loc,lin}}\). % lean: canonicalLocLinWeights

Fix a block \(V_b\) and an admissible choice \((\varepsilon_i,S_i)_{i\in V_b}\) in \(\Psi_{b,t}\).  All units of \(V_b\) carry the same weight \(g_{d}(Z_b)\), so
\[
\sum_{i\in V_b}\varepsilon_i\bigl\{w^0_{i,t}(Z_b)Z_{S_i}-\mathbf 1(S_i\ne\varnothing)\bigr\}
=
g_{d}(Z_b)F(Z_b)-\Lambda,
\]
\[
F(Z_b):=\sum_{i\in V_b}\varepsilon_iZ_{S_i},
\qquad
\Lambda:=\sum_{i\in V_b}\varepsilon_i\mathbf 1(S_i\ne\varnothing).
\]
The block-score identities give \(\mathbb E_Z[g_{d}(Z_b)F(Z_b)]=\Lambda\), whence
\[
\mathbb E_Z\!\left[\bigl(g_{d}(Z_b)F(Z_b)-\Lambda\bigr)^2\right]
=
\mathbb E_Z\!\left[\bigl(g_{d}(Z_b)F(Z_b)\bigr)^2\right]-\Lambda^2
\le
\mathbb E_Z\!\left[\bigl(g_{d}(Z_b)F(Z_b)\bigr)^2\right].
\]
Since \(0\le Z_{S_i}\le1\), \(|\varepsilon_i|=1\) and \(|V_b|=d\), we have \(|F(Z_b)|\le d\) pointwise, so
\[
\mathbb E_Z\!\left[\bigl(g_{d}(Z_b)F(Z_b)\bigr)^2\right]
\le
d^2\,\mathbb E_Z[g_{d}(Z_b)^2]
=
d^2A .
\]
Maximizing over admissible choices gives the upper bound, and the empty-monomial construction gives the reverse inequality for the same canonical weights, so
\[
\Psi_{b,t}(w^0)\le d^2A,
\qquad\text{and}\qquad
\Psi_{b,t}(w^0)=d^2A .
\]
% lean: canonical_completeBlock_blockExtremal_upper
By the worst-case identity and the per-block lower bound, every \(w\in\mathcal E_t^{\mathrm{loc,lin}}\) satisfies
\[
\sup_{c\in\mathcal M_t(G_t)}\mathbb E_Z\!\left[(\widehat\tau_{w,t}-\tau_{n})^2\right]
=
\frac{B^2}{n^2}\sum_{b=1}^{m}\Psi_{b,t}(w)
\ge
\frac{B^2}{n^2}m d^2A
=
\frac{B^2A}{m},
\]
the last equality because \(n=md\); and the same computation with \(\Psi_{b,t}(w^0)=d^2A\) turns the inequality into an equality for \(w^0\).  Taking the infimum over \(\mathcal E_t^{\mathrm{loc,lin}}\),
\[
R_t^{\mathrm{loc,lin}}=\frac{B^2A_{d_t}}{m_t}
=\frac{B^2d_tA_{d_t}}{n_t},
\]
where the second equality is again \(n_t=m_td_t\). % lean: completeBlock_benchmark_algebra, locLinMinimaxRisk_exact

\item \textbf{The risk ratio equals one plus the normalized block excess.}  For a sequence \(w_t\in\mathcal E_t^{\mathrm{loc,lin}}\) set
\[
Y_t:=
\frac{1}{m_td_t^2A_{d_t}}
\sum_{b=1}^{m_t}\bigl\{\Psi_{b,t}(w_t)-d_t^2A_{d_t}\bigr\}.
\]
Dividing the worst-case identity by \(B^2A_{d_t}/m_t\) and using \(n_t^2=m_t^2d_t^2\),
\[
\frac{
\sup_{c\in\mathcal M_t(G_t)}\mathbb E_Z[(\widehat\tau_{w_t,t}-\tau_{n_t})^2]
}{
B^2A_{d_t}/m_t
}
=
\frac{m_t}{n_t^2A_{d_t}}\sum_{b=1}^{m_t}\Psi_{b,t}(w_t)
=
\frac{1}{m_td_t^2A_{d_t}}\sum_{b=1}^{m_t}\Psi_{b,t}(w_t)
=
1+Y_t .
\]
% lean: locLinRiskRatio_eq_one_add_normalizedBlockExcess
Adding, respectively subtracting, the constant sequence \(1\) therefore shows that the left-hand ratio converges to \(1\) if and only if \(Y_t\to0\), which is the asserted equivalence. % lean: riskRatio_tendsto_iff_normalizedBlockExcess

\item \textbf{Score closeness implies attainment.}  Put \(\delta_{i,t}:=w_{i,t}-g_{d_t}\) on the block of \(i\), and set
\[
E_{b,t}:=\sum_{i\in V_b}\mathbb E_Z\!\left[\delta_{i,t}(Z_b)^2\right],
\qquad
E_t:=\sum_{b=1}^{m_t}E_{b,t}
=\sum_{i\in V_{n_t}}\mathbb E_Z\!\left[\bigl\{w_{i,t}(Z_{b(i)})-g_{d_t}(Z_{b(i)})\bigr\}^2\right],
\]
\[
X_t:=\frac{E_t}{n_tA_{d_t}}
=
\frac{1}{n_tA_{d_t}}
\sum_{i\in V_{n_t}}
\mathbb E_Z\!\left[\bigl\{w_{i,t}(Z_{b(i)})-g_{d_t}(Z_{b(i)})\bigr\}^2\right],
\]
so that the hypothesis of this clause is \(X_t\to0\).  The per-block lower bound gives \(Y_t\ge0\) for every \(t\).

Fix \(t\), a block \(V_b\) and an admissible choice \((\varepsilon_i,S_i)_{i\in V_b}\).  Splitting each weight into its canonical part and its deviation,
\[
\sum_{i\in V_b}\varepsilon_i\bigl\{w_{i,t}(Z_b)Z_{S_i}-\mathbf 1(S_i\ne\varnothing)\bigr\}
=
f(Z_b)+v(Z_b),
\]
\[
f(Z_b):=\sum_{i\in V_b}\varepsilon_i\bigl\{g_{d_t}(Z_b)Z_{S_i}-\mathbf 1(S_i\ne\varnothing)\bigr\},
\qquad
v(Z_b):=\sum_{i\in V_b}\varepsilon_i\delta_{i,t}(Z_b)Z_{S_i}.
\]
For the canonical part, the calculation with the weights \(w^0\) gives \(\mathbb E_Z[f(Z_b)^2]\le d_t^2A_{d_t}\).  For \(v\), the Cauchy--Schwarz inequality over the \(d_t\) summands together with \(0\le Z_{S_i}\le1\) and \(|\varepsilon_i|=1\) gives \(v(Z_b)^2\le d_t\sum_{i\in V_b}\delta_{i,t}(Z_b)^2\) pointwise.  Thus
\[
\mathbb E_Z[f(Z_b)^2]\le d_t^2A_{d_t},
\qquad
\mathbb E_Z[v(Z_b)^2]\le d_tE_{b,t}.
\]
% lean: blockCandidatePerturb_energy_le
Bounding the cross term by Cauchy--Schwarz, \(\mathbb E_Z[(f+v)^2]\le\mathbb E_Z[f^2]+2\sqrt{\mathbb E_Z[f^2]}\sqrt{\mathbb E_Z[v^2]}+\mathbb E_Z[v^2]\), and maximizing over admissible choices,
\[
\Psi_{b,t}(w_t)-d_t^2A_{d_t}
\le
2\sqrt{d_t^2A_{d_t}}\sqrt{d_tE_{b,t}}+d_tE_{b,t}.
\]
% lean: finiteDesign_E_add_sq_le, blockExtremal_le_of_weight_distance
Summing over \(b\) and applying Cauchy--Schwarz once more, \(\sum_{b=1}^{m_t}\sqrt{d_tE_{b,t}}\le\sqrt{m_t}\sqrt{d_tE_t}\), so
\[
\sum_{b=1}^{m_t}\bigl\{\Psi_{b,t}(w_t)-d_t^2A_{d_t}\bigr\}
\le
2\sqrt{d_t^2A_{d_t}}\sqrt{m_t}\sqrt{d_tE_t}+d_tE_t .
\]
% lean: completeBlockExtremal_excess_le_distance
Dividing by \(m_td_t^2A_{d_t}\) and using \(n_t=m_td_t\),
\[
\frac{
2\sqrt{d_t^2A_{d_t}}\sqrt{m_t}\sqrt{d_tE_t}+d_tE_t
}{
m_td_t^2A_{d_t}
}
=
2\sqrt{\frac{E_t}{n_tA_{d_t}}}+\frac{E_t}{n_tA_{d_t}}
=
2\sqrt{X_t}+X_t ,
\]
% lean: normalized_sqrt_excess_identity
so that
\[
0\le Y_t\le 2\sqrt{X_t}+X_t .
\]
If \(X_t\to0\) then the right-hand bound tends to \(0\) by continuity of the square root, so \(Y_t\to0\) by squeezing.  Since the risk ratio is \(1+Y_t\), this is the asserted convergence of the risk ratio to \(1\). % lean: normalizedWeightDistance_tendsto_implies_riskRatio

\item \textbf{An optimal sequence at normalized distance two.}  Assume \(d_t\to\infty\).  For every \(t\) and each block \(V_b\), fix a distinguished unit \(i_b^\star\in V_b\), and define the perturbation by
\[
\psi_b(Z_b)
:=
\begin{cases}
\left\{\dfrac{2d_tA_{d_t}}{[p(1-p)]^{d_t}}\right\}^{1/2}
\prod_{j\in V_b}(Z_j-p), & \text{if }\beta<d_t,\\[1ex]
0, & \text{if }\beta\ge d_t.
\end{cases}
\]
Set, for every \(t\),
\[
w_{i,t}(Z)
:=
g_{d_t}(Z_{b(i)})
+
\mathbf 1\bigl(i=i^\star_{b(i)}\bigr)\,\psi_{b(i)}(Z_{b(i)}).
\]
Thus the sequence is defined on the whole index set; at indices with \(\beta\ge d_t\) it coincides with the canonical block-score weights. % lean: orthogonalBlockPerturb, orthogonalWitnessWeights

\textbf{Admissibility.}  Each \(\psi_b\) is a function of \(Z_b\) alone.  If \(\beta\ge d_t\), then \(\psi_b=0\), and the block-score identities verify the centering and unit raw-moment conditions.  If \(\beta<d_t\), then \(\bar\beta_{d_t}=\beta\), and for every \(S\subseteq V_b\) with \(|S|\le\beta\), including \(S=\varnothing\), at least one index \(j\in V_b\) outside \(S\) occurs in the product.  Under \cref{obj:ass:bernoulli-design}, that coordinate is independent of \(Z_S\) and of the remaining factors and satisfies \(\mathbb E_Z[Z_j-p]=0\); hence
\[
\mathbb E_Z\!\left[\bigl(\prod_{j\in V_b}(Z_j-p)\bigr)Z_S\right]=0
\qquad\text{for every }S\subseteq V_b\text{ with }|S|<d_t .
\]
% lean: E_global_centered_mul_raw_eq_zero
So the perturbation changes neither the centering nor the unit raw moments inherited from \(g_{d_t}\), and \(w_t\in\mathcal E_t^{\mathrm{loc,lin}}\) for every \(t\).

Since \(d_t\to\infty\), the inequality \(\beta<d_t\) holds for all sufficiently large \(t\).  All remaining computations are on this tail, where the first branch in the definition of \(\psi_b\) applies.

\textbf{Perturbation energy.}  By independence and \(\mathbb E_Z[(Z_j-p)^2]=p(1-p)\),
\[
\mathbb E_Z\!\left[\bigl(\prod_{j\in V_b}(Z_j-p)\bigr)^2\right]=[p(1-p)]^{d_t},
\qquad\text{hence}\qquad
\mathbb E_Z\!\left[\psi_b(Z_b)^2\right]=2d_tA_{d_t}.
\]
% lean: E_global_centeredMonomial_mul, orthogonalBlockPerturb_energy

\textbf{Block bound and the risk ratio.}  Fix a block \(V_b\) and an admissible choice, and split the inner sum into canonical and perturbation parts.  Because the perturbation is carried by the single unit \(i_b^\star\), the deviation part reduces to one summand, \(v(Z_b)=\varepsilon_{i_b^\star}\psi_b(Z_b)Z_{S_{i_b^\star}}\), and \(0\le Z_{S_{i_b^\star}}\le1\) gives \(\mathbb E_Z[v(Z_b)^2]\le\mathbb E_Z[\psi_b(Z_b)^2]=2d_tA_{d_t}\); the canonical part obeys \(\mathbb E_Z[f(Z_b)^2]\le d_t^2A_{d_t}\).  The same Cauchy--Schwarz bound on the cross term therefore yields, for all sufficiently large \(t\) and every \(b\),
\[
\Psi_{b,t}(w_t)
\le
d_t^2A_{d_t}
+2\sqrt{d_t^2A_{d_t}}\sqrt{2d_tA_{d_t}}
+2d_tA_{d_t}.
\]
% lean: monomial_mul_sq_energy_le, orthogonalWitness_blockExtremal_upper
Dividing the excess by \(d_t^2A_{d_t}\) gives \(2\sqrt{2/d_t}+2/d_t\) for each block, hence the same bound for \(Y_t\), while \(Y_t\ge0\) by the per-block lower bound.  Since the risk ratio is \(1+Y_t\),
\[
1
\le
\frac{
\sup_{c\in\mathcal M_t(G_t)}\mathbb E_Z[(\widehat\tau_{w_t,t}-\tau_{n_t})^2]
}{
B^2A_{d_t}/m_t
}
\le
1+2\sqrt{\frac{2}{d_t}}+\frac{2}{d_t}
\]
for all sufficiently large \(t\).  Since \(d_t\to\infty\), squeezing proves the asserted convergence of the risk ratio to \(1\). % lean: orthogonalWitness_riskRatio_tendsto_one

\textbf{Normalized distance.}  For the same sufficiently large \(t\), the deviation \(w_{i,t}-g_{d_t}\) equals \(\psi_b\) at \(i=i_b^\star\) and vanishes at every other unit, so
\[
\sum_{i\in V_{n_t}}
\mathbb E_Z\!\left[\bigl\{w_{i,t}(Z_{b(i)})-g_{d_t}(Z_{b(i)})\bigr\}^2\right]
=
\sum_{b=1}^{m_t}\mathbb E_Z\!\left[\psi_b(Z_b)^2\right]
=
2m_td_tA_{d_t},
\]
and dividing by \(n_tA_{d_t}=m_td_tA_{d_t}\) gives the value \(2\). % lean: orthogonalWitnessWeights_distance_energy, orthogonalWitness_normalized_distance_eq_two

\textbf{Invariance under block-preserving relabeling.}  Let \(\pi\) be a bijection of \(V_{n_t}\) with \(b(\pi(i))=b(i)\) for every \(i\), and take \(t\) in the same sufficiently large tail.  Then \(\pi\) maps each block onto itself and therefore permutes the coordinates of \(Z_b\) among themselves.  The relabeled distance moves the score term only: at the unit \(i\) it compares the weight \(w_{i,t}\), read at the unit \(i\) and at the original assignment, with the canonical score of the relabeled unit \(\pi(i)\), whose neighborhood is again \(V_{b(i)}\), read at the relabeled assignment whose \(j\)th coordinate is \(Z_{\pi^{-1}(j)}\).  Since \(g_{d_t}\) is a symmetric function of the coordinates of its block, that relabeled score is the original one,
\[
g_{d_t}\bigl((Z_{\pi^{-1}(j)})_{j\in V_{b(i)}}\bigr)
=
g_{d_t}(Z_{b(i)})
\qquad\text{for every }i\in V_{n_t}.
\]
% lean: canonicalBlockScore_relabel_within_blocks
Each summand of the normalized distance is therefore unchanged, and substituting this equality into the normalized-distance expression gives
\[
\frac{1}{n_tA_{d_t}}
\sum_{i\in V_{n_t}}
\mathbb E_Z\!\left[
\Bigl\{w_{i,t}(Z_{b(i)})-g_{d_t}\bigl((Z_{\pi^{-1}(j)})_{j\in V_{b(i)}}\bigr)\Bigr\}^2
\right]
=2 ,
\]
which is the asserted invariance.
% lean: orthogonalWitness_normalized_relabel_distance_eq_two
\end{enumerate}
\end{proof}
